% =====================================================================
%  Tarea 1 — IIC3533 Computación de Alto Rendimiento — 2026-2
%  Compilar con: pdflatex informe.tex  (dos veces, por las referencias)
%  Las figuras se leen desde ../resultados/<maquina>/
% =====================================================================
\documentclass[11pt,letterpaper]{article}

\usepackage[utf8]{inputenc}
\usepackage[T1]{fontenc}
% Español si está disponible (Overleaf lo trae; en Debian/Ubuntu es el paquete
% texlive-lang-spanish). Si falta, el documento igual compila, solo que los
% títulos automáticos ("Índice", "Cuadro") quedan en inglés.
\IfFileExists{spanish.ldf}{\usepackage[spanish,es-tabla]{babel}}{%
  \usepackage[english]{babel}%
  \PackageWarningNoLine{informe}{babel-spanish no disponible: instalar
    texlive-lang-spanish para los titulos en espanol}}
\usepackage{amsmath,amssymb}
\usepackage{graphicx}
\usepackage{booktabs}
\usepackage{listings}
\usepackage{xcolor}
\usepackage{caption}
\usepackage{float}
\usepackage[margin=2.5cm]{geometry}
\usepackage[hidelinks]{hyperref}

% ---------------------------------------------------------------
% Marcador de secciones pendientes. BORRAR TODOS ANTES DE ENTREGAR.
% ---------------------------------------------------------------
\newcommand{\pendiente}[1]{%
  \par\medskip\noindent
  \fcolorbox{red!60!black}{red!5}{%
    \begin{minipage}{0.96\linewidth}
      \small\textbf{\textcolor{red!60!black}{PENDIENTE}}\par\vspace{2pt}#1
    \end{minipage}}%
  \par\medskip}

% Figura que no falla si el archivo todavía no existe
% \detokenize evita que los guiones bajos de las rutas se interpreten como
% subíndices matemáticos cuando la figura todavía no existe.
\newcommand{\figurasiexiste}[3]{%
  \IfFileExists{#1}{\includegraphics[width=#2]{#1}}{%
    \fcolorbox{red!60!black}{red!5}{\begin{minipage}{0.9\linewidth}
      \centering\small\vspace{1cm}
      \textcolor{red!60!black}{\textbf{Figura pendiente:} \texttt{\detokenize{#1}}}\par
      \vspace{2pt}#3\vspace{1cm}
    \end{minipage}}}}

% ---------------------------------------------------------------
% Estilo para el código Python
% ---------------------------------------------------------------
\definecolor{codegris}{rgb}{0.95,0.95,0.95}
\definecolor{codeverde}{rgb}{0,0.5,0}
\definecolor{codemorado}{rgb}{0.5,0,0.5}
\lstset{
  language=Python,
  basicstyle=\ttfamily\scriptsize,
  keywordstyle=\color{blue},
  commentstyle=\color{codeverde},
  stringstyle=\color{codemorado},
  numbers=left, numberstyle=\tiny\color{gray}, numbersep=6pt,
  backgroundcolor=\color{codegris},
  frame=single, rulecolor=\color{gray!50},
  breaklines=true, breakatwhitespace=true,
  showstringspaces=false, tabsize=4,
  captionpos=b, xleftmargin=12pt,
  literate={á}{{\'a}}1 {é}{{\'e}}1 {í}{{\'i}}1 {ó}{{\'o}}1 {ú}{{\'u}}1
           {ñ}{{\~n}}1 {Á}{{\'A}}1 {É}{{\'E}}1 {Í}{{\'I}}1 {Ó}{{\'O}}1
           {Ú}{{\'U}}1 {Ñ}{{\~N}}1 {¿}{{?`}}1 {¡}{{!`}}1
}

% ---------------------------------------------------------------
\title{\vspace{-1cm}\textbf{Tarea 1: Bootstrapping paralelo para regresión lineal}\\
       \large IIC3533 --- Computación de Alto Rendimiento --- 2026-2}
\author{Benjamín Saldías \and Bruno \textit{[apellido]} \and \textit{[tercer integrante]}}
\date{\today}

\begin{document}
\maketitle

% =====================================================================
\pendiente{\textbf{Estado del documento (borrar esta caja antes de entregar).}
Los ítems (a) a (i) están contestados con los resultados de la \textbf{máquina 1}.
Falta únicamente la tanda de la \textbf{máquina 2}, que completa las tablas y figuras
marcadas y habilita el ítem (j). Para generarla: \texttt{bash correr\_todo.sh} en el
segundo computador, y luego reemplazar \texttt{MAQUINA2} por su etiqueta en las rutas
de las figuras.}

\section*{Resumen}

Se implementaron y compararon tres versiones del algoritmo de \textit{bootstrapping}
para estimar intervalos de confianza de los coeficientes de una regresión lineal por
mínimos cuadrados, con $N = 100\,000$ observaciones, $k = 300$ variables y $B = 48$
resamples. Las tres versiones producen intervalos estadísticamente equivalentes y
consistentes con los coeficientes verdaderos $\boldsymbol{\beta}^*$, pero difieren en
un factor cercano a $6\times$ en tiempo de ejecución: la versión que resuelve la
ecuación normal directamente con NumPy resulta claramente la más eficiente frente a
las que usan \texttt{LinearRegression} de \texttt{scikit-learn}. El análisis de
escalamiento muestra que el problema está limitado por ancho de banda de memoria y no
por cómputo, lo que explica que el \textit{speedup} se sature muy por debajo del ideal.

\tableofcontents
\newpage

% =====================================================================
\section{Entorno experimental}
\label{sec:entorno}

Los experimentos se realizaron en dos computadores distintos, cuyas características se
resumen en la Tabla~\ref{tab:maquinas}. Se mantuvieron deliberadamente las diferencias
de entorno (biblioteca BLAS y versión de Python), ya que son parte del objeto de
comparación del ítem (j).

\begin{table}[H]
\centering
\caption{Características de los dos computadores utilizados.}
\label{tab:maquinas}
\small
\begin{tabular}{lll}
\toprule
 & \textbf{Máquina 1} & \textbf{Máquina 2} \\
\midrule
Identificador      & \texttt{DESKTOP-52S4V2O} & \textit{[completar]} \\
CPU                & Intel Core i5-1135G7 @ 2.40\,GHz & \textit{[completar]} \\
Cores físicos      & 4 & \textit{[completar]} \\
\textbf{Cores lógicos ($p_{\max}$)} & \textbf{8} & \textit{[completar]} \\
Memoria RAM        & 7{,}6 GiB & \textit{[completar]} \\
Sistema operativo  & WSL2 sobre Windows & \textit{[completar]} \\
Python             & 3.13.15 (conda) & \textit{[completar]} \\
NumPy              & 2.5.2 & \textit{[completar]} \\
joblib             & 1.5.3 & \textit{[completar]} \\
threadpoolctl      & 3.5.0 & \textit{[completar]} \\
\textbf{BLAS}      & \textbf{MKL} 2025.0 & \textit{[completar]} \\
Threads BLAS por defecto & 4 & \textit{[completar]} \\
\bottomrule
\end{tabular}
\end{table}

\pendiente{Completar la columna de la máquina 2 con la salida de
\texttt{python src/config.py} y \texttt{lscpu}, que queda registrada automáticamente al
comienzo de \texttt{resultados/<máquina>/terminal.txt}.}

\paragraph{Estructura del código.} Cada versión del algoritmo vive en su propio archivo
(\texttt{bs\_auto.py}, \texttt{bs\_sklearn.py}, \texttt{bs\_numpy.py}) y comparte la
generación de datos (\texttt{data\_utils.py}). Los experimentos se lanzan con
\texttt{correr\_todo.sh}, que ejecuta la tanda completa y deja el registro en
\texttt{resultados/<máquina>/}, de modo que los resultados de un computador no
sobrescriben los del otro.

% =====================================================================
\section{(a) Generación de los datos sintéticos}
\label{sec:a}

Los datos se generan con una semilla fija (\texttt{seed = 42}), siguiendo exactamente
los tres pasos del enunciado: se muestrean los $k+1$ coeficientes verdaderos
$\boldsymbol{\beta}^*$ desde $\mathcal{N}(0,1)$, se construye la matriz de datos de
$N \times k$ con entradas i.i.d.\ $\mathcal{N}(0,1)$, se le antepone una columna de
unos para formar $X \in \mathbb{R}^{N \times (k+1)}$, y finalmente
$\mathbf{y} = X\boldsymbol{\beta}^* + \boldsymbol{\varepsilon}$ con
$\boldsymbol{\varepsilon}$ de $N$ valores i.i.d.\ $\mathcal{N}(0,1)$.

\begin{lstlisting}[caption={Generación de los datos sintéticos (\texttt{data\_utils.py}).}]
def generate_synthetic_data(N=100000, k=300, seed=42):
    np.random.seed(seed)                          # reproducibilidad

    beta_star = np.random.normal(0, 1, k + 1)     # (I)  coeficientes verdaderos
    X_features = np.random.normal(0, 1, (N, k))   # (II) matriz de datos

    unos = np.ones((N, 1))
    X = np.hstack((unos, X_features))             # columna de unos al inicio

    ruido = np.random.normal(0, 1, N)             # (III)
    y = X.dot(beta_star) + ruido

    return X, y, beta_star
\end{lstlisting}

\paragraph{Verificación.} Las dimensiones resultantes son
$X \in \mathbb{R}^{100000 \times 301}$, $\mathbf{y} \in \mathbb{R}^{100000}$ y
$\boldsymbol{\beta}^* \in \mathbb{R}^{301}$, como exige el enunciado.

\paragraph{Consideración de memoria.} La matriz $X$ ocupa
$100\,000 \times 301 \times 8~\text{bytes} \approx 241$~MB en doble precisión. Este
número resulta central para todo el análisis posterior: es el volumen de datos que hay
que hacer llegar a cada proceso, y cada resample crea una copia de ese mismo tamaño.
Generar los datos toma apenas $0{,}55$~s, de modo que su costo es despreciable frente
al de los ajustes.

% =====================================================================
\section{(b) Las tres implementaciones}
\label{sec:b}

\subsection{Versión 1: \texttt{bs\_auto.py} --- \texttt{BaggingRegressor}}

El bootstrapping y el paralelismo quedan enteramente delegados a
\texttt{scikit-learn}. Se usa \texttt{fit\_intercept=False} porque $X$ ya incluye la
columna de unos, y \texttt{max\_samples=1.0} con \texttt{bootstrap=True} para que cada
estimador muestree $N$ filas con reemplazo.

\begin{lstlisting}[caption={Núcleo de \texttt{bs\_auto.py}.}]
modelo_bagging = BaggingRegressor(
    estimator=LinearRegression(fit_intercept=False),
    n_estimators=B,      # B = 48 resamples
    n_jobs=p,            # p procesos
    bootstrap=True,
    max_samples=1.0,     # muestrea N elementos con reemplazo
    random_state=42,     # imprescindible para reproducibilidad (ver seccion 4)
)
modelo_bagging.fit(X, y)
beta_hats = np.array([est.coef_ for est in modelo_bagging.estimators_])
\end{lstlisting}

\subsection{Versión 2: \texttt{bs\_sklearn.py} --- \texttt{joblib} + \texttt{LinearRegression}}

Aquí el paralelismo es explícito. Cada tarea sortea sus propios índices y ajusta un
\texttt{LinearRegression}. La semilla $42+b$ hace que cada resample sea distinto pero
reproducible.

\begin{lstlisting}[caption={Núcleo de \texttt{bs\_sklearn.py}.}]
def ajustar_bootstrap(X, y, semilla, t=1):
    with threadpool_limits(limits=t, user_api="blas"):
        np.random.seed(semilla)
        N = X.shape[0]
        indices = np.random.choice(N, size=N, replace=True)  # (I)
        X_b, y_b = X[indices], y[indices]                    # (II)
        modelo = LinearRegression(fit_intercept=False)       # (III)
        modelo.fit(X_b, y_b)
        return modelo.coef_

beta_hats = Parallel(n_jobs=p)(
    delayed(ajustar_bootstrap)(X, y, 42 + b, t) for b in range(B))
\end{lstlisting}

\subsection{Versión 3: \texttt{bs\_numpy.py} --- \texttt{joblib} + ecuación normal}

La única diferencia con la anterior es el ajuste: en lugar de delegar en
\texttt{scikit-learn}, se resuelve directamente el sistema
$(X^\top X)\hat{\boldsymbol{\beta}} = X^\top \mathbf{y}$.

\begin{lstlisting}[caption={Núcleo de \texttt{bs\_numpy.py}.}]
def ajustar_bootstrap_numpy(X, y, semilla, t=1):
    with threadpool_limits(limits=t, user_api="blas"):
        np.random.seed(semilla)
        N = X.shape[0]
        indices = np.random.choice(N, size=N, replace=True)
        X_b, y_b = X[indices], y[indices]
        X_T = X_b.T
        return np.linalg.solve(X_T.dot(X_b), X_T.dot(y_b))
\end{lstlisting}

\subsection{Iteraciones de optimización}
\label{sec:b-iteraciones}

\paragraph{Iteración 1: de \texttt{lstsq} a la ecuación normal.}
La primera implementación usaba \texttt{LinearRegression}, que internamente llama a
\texttt{scipy.linalg.lstsq} y resuelve el problema mediante una descomposición SVD de
la matriz completa $X^{(b)} \in \mathbb{R}^{100000 \times 301}$. Ese método es
numéricamente más estable ante matrices mal condicionadas, pero mucho más caro:
descompone una matriz de $100\,000$ filas.

Resolver la ecuación normal, en cambio, primero \emph{reduce} el problema: el producto
$X^\top X$ cuesta $O(Nk^2)$ pero produce una matriz de apenas $301 \times 301$, y
factorizarla es entonces trivial ($O(k^3)$ con $k=301$). Como los datos del problema
son i.i.d.\ $\mathcal{N}(0,1)$, la matriz está muy bien condicionada y no hay razón
para pagar el costo de la SVD.

El efecto medido en la máquina 1 fue de aproximadamente $6\times$:

\begin{table}[H]
\centering
\caption{Efecto de la iteración 1, medido en la máquina 1 con $p = 2$.}
\label{tab:iteracion1}
\small
\begin{tabular}{lrr}
\toprule
\textbf{Versión} & \textbf{Método de ajuste} & \textbf{Tiempo (s)} \\
\midrule
\texttt{bs\_auto}    & \texttt{lstsq} (SVD), vía \texttt{BaggingRegressor} & 96{,}40 \\
\texttt{bs\_sklearn} & \texttt{lstsq} (SVD), vía \texttt{LinearRegression} & 86{,}83 \\
\texttt{bs\_numpy}   & ecuación normal (\texttt{np.linalg.solve})          & \textbf{10{,}88} \\
\bottomrule
\end{tabular}
\end{table}

\paragraph{Iteración 2: evitar materializar $X^{(b)}$.}
Una vez eliminado el costo de la SVD, el cuello de botella deja de ser aritmético. La
línea \texttt{X\_b = X[indices]} copia 241~MB en cada resample con un patrón de
\emph{acceso aleatorio} a memoria, que es el caso más desfavorable para la jerarquía de
caché.

Ese paso se puede evitar por completo. Un resample bootstrap no es más que asignar a
cada fila un peso entero $w_i$ igual al número de veces que fue sorteada, de modo que

\begin{equation}
X^{(b)\top} X^{(b)} \;=\; X^\top \operatorname{diag}(\mathbf{w})\, X,
\qquad
X^{(b)\top} \mathbf{y}^{(b)} \;=\; X^\top (\mathbf{w} \odot \mathbf{y}).
\label{eq:pesos}
\end{equation}

\begin{lstlisting}[caption={Variante con pesos multinomiales (\texttt{bs\_numpy.py --pesos}).}]
def ajustar_bootstrap_numpy_pesos(X, y, semilla, t=1):
    with threadpool_limits(limits=t, user_api="blas"):
        np.random.seed(semilla)
        N = X.shape[0]
        indices = np.random.choice(N, size=N, replace=True)
        w = np.bincount(indices, minlength=N).astype(X.dtype)  # conteos por fila
        Xw = X * w[:, None]          # acceso secuencial, no aleatorio
        return np.linalg.solve(X.T.dot(Xw), X.T.dot(y * w))
\end{lstlisting}

Como se usan \emph{los mismos índices} que en la versión anterior, ambas calculan el
mismo estimador: la diferencia numérica medida es de $6{,}2 \times 10^{-15}$, atribuible
únicamente al orden de acumulación en punto flotante. Sigue existiendo un arreglo
temporal del tamaño de $X$, pero se recorre de forma secuencial.

Medida a escala completa en la máquina 1 con $p = 8$, la variante con pesos resultó
\textbf{más lenta}, no más rápida:

\begin{table}[H]
\centering
\caption{Efecto de la iteración 2, medido en la máquina 1 con $p = 8$, $t = 1$.}
\label{tab:iteracion2}
\small
\begin{tabular}{lr}
\toprule
\textbf{Variante de \texttt{bs\_numpy}} & \textbf{Tiempo (s)} \\
\midrule
Índices: \texttt{X\_b = X[indices]}                     & \textbf{10{,}98} \\
Pesos: \texttt{Xw = X * w[:, None]} (ecuación~\ref{eq:pesos}) & 13{,}57 \\
\bottomrule
\end{tabular}
\end{table}

El resultado es negativo pero informativo, y refuerza la interpretación que se desarrolla
en la Sección~\ref{sec:h}. Si el cuello de botella fuera el \emph{patrón} de acceso a
memoria, convertir el acceso aleatorio en secuencial debería haber ayudado. Que no lo
haga indica que lo que domina es el \emph{volumen} de datos movido: ambas variantes
asignan y recorren un arreglo temporal de 241~MB por resample, y la versión con pesos
agrega además el cálculo de \texttt{np.bincount} y una multiplicación elemento a elemento
sobre los $100\,000 \times 301$ valores. El acceso aleatorio de \texttt{X[indices]} no era
tan costoso como esperábamos, probablemente porque el \textit{prefetcher} del procesador
tolera bien un patrón que, aunque desordenado, recorre filas completas y contiguas de
2408 bytes cada una.

Conclusión de las dos iteraciones: la mejora relevante fue algorítmica (cambiar SVD por
ecuación normal, $6\times$), mientras que la micro-optimización del patrón de acceso no
aportó. Conservamos ambas variantes en el código porque la comparación es parte del
resultado.

\paragraph{Observación adicional: consumo de memoria.}
Al ejecutar en la máquina 2 (limitada a 6~GB), \texttt{bs\_sklearn} fue terminada por
el sistema operativo (\texttt{SIGKILL}) con $p = 6$, mientras que \texttt{bs\_numpy}
completó sin problemas con el mismo $p$. La razón es que \texttt{lstsq} requiere, además
de la copia de $X^{(b)}$, espacio de trabajo interno para la descomposición SVD, de
varias veces el tamaño de la matriz. La ventaja de la ecuación normal, por lo tanto, no
es solo de tiempo sino también de huella de memoria.

% =====================================================================
\section{(c) Correctitud y reproducibilidad}
\label{sec:c}

Para responder este ítem se implementó el script \texttt{src/verificar\_correctitud.py},
que ejecuta las versiones sobre los mismos datos y compara sus salidas.

\subsection{Consistencia con los coeficientes verdaderos}

Los intervalos de confianza se construyen según el Paso 3 del enunciado: para cada
coeficiente $j$ se ordenan los $B$ valores $\hat{\beta}^{(1)}_j, \dots, \hat{\beta}^{(B)}_j$
y se descarta el $2{,}5\%$ inferior y superior.

\begin{lstlisting}[caption={Construcción del intervalo de confianza al 95\%.}]
limite_inferior = np.percentile(beta_hats, 2.5, axis=0)
limite_superior = np.percentile(beta_hats, 97.5, axis=0)
\end{lstlisting}

Los resultados a escala completa ($N = 100\,000$, $k = 300$, $B = 48$, con $p = 8$) se
presentan en la Tabla~\ref{tab:cobertura}.

\begin{table}[H]
\centering
\caption{Intervalos de confianza al 95\,\% y cobertura de $\boldsymbol{\beta}^*$.}
\label{tab:cobertura}
\small
\begin{tabular}{lccc}
\toprule
\textbf{Versión} & \textbf{IC 95\,\% de $\beta_0$} & \textbf{Ancho medio} & \textbf{Cobertura de $\boldsymbol{\beta}^*$} \\
\midrule
\texttt{bs\_numpy}          & $[0{,}4925,\ 0{,}5043]$ & $0{,}01142$ & $91{,}4\%$ \\
\texttt{bs\_sklearn}        & $[0{,}4925,\ 0{,}5043]$ & $0{,}01142$ & $91{,}4\%$ \\
\texttt{bs\_numpy (pesos)}  & $[0{,}4925,\ 0{,}5043]$ & $0{,}01142$ & $91{,}4\%$ \\
\texttt{bs\_auto}           & $[0{,}4902,\ 0{,}5018]$ & $0{,}01141$ & $91{,}7\%$ \\
\bottomrule
\end{tabular}
\par\vspace{2pt}
\footnotesize Valor verdadero: $\beta_0^* = 0{,}4967$. Los cuatro intervalos lo contienen.
\end{table}

Los cuatro intervalos contienen el valor verdadero $\beta_0^* = 0{,}4967$, y los anchos
medios coinciden hasta la quinta cifra decimal, lo que confirma que las cuatro
implementaciones estiman la misma distribución muestral. A modo de ilustración, la
Tabla~\ref{tab:primeros} muestra los primeros cinco coeficientes.

\begin{table}[H]
\centering
\caption{Intervalos de \texttt{bs\_numpy} para los primeros cinco coeficientes.}
\label{tab:primeros}
\small
\begin{tabular}{crrrc}
\toprule
$j$ & $\beta^*_j$ & \textbf{Límite inferior} & \textbf{Límite superior} & \textbf{¿Cubre?} \\
\midrule
0 & $ 0{,}4967$ & $ 0{,}4925$ & $ 0{,}5043$ & sí \\
1 & $-0{,}1383$ & $-0{,}1400$ & $-0{,}1286$ & sí \\
2 & $ 0{,}6477$ & $ 0{,}6398$ & $ 0{,}6524$ & sí \\
3 & $ 1{,}5230$ & $ 1{,}5206$ & $ 1{,}5305$ & sí \\
4 & $-0{,}2342$ & $-0{,}2393$ & $-0{,}2290$ & sí \\
\bottomrule
\end{tabular}
\end{table}

Con $B = 48$ resamples, los percentiles $2{,}5$ y $97{,}5$ se estiman a partir de muy
pocas observaciones en cada cola, por lo que la cobertura empírica fluctúa en torno al
95\,\% nominal y no debe esperarse que coincida exactamente. Aumentar $B$ estrecharía
esa fluctuación, a costa de más cómputo.

\subsection{¿En qué medida son equivalentes las distintas versiones?}
\label{sec:c-equivalencia}

Conviene distinguir dos niveles de equivalencia, porque las versiones no se comportan
igual:

\begin{description}
\item[Equivalencia numérica.] \texttt{bs\_sklearn} y \texttt{bs\_numpy} comparten la
  semilla $42+b$ en cada tarea, de modo que remuestrean \emph{exactamente las mismas
  filas}. Sus estimadores deben coincidir hasta el error de redondeo, y eso es lo que
  se observa. Lo mismo vale para la variante con pesos, que reutiliza los mismos
  índices.
\item[Equivalencia estadística.] \texttt{bs\_auto} no puede coincidir numéricamente:
  \texttt{BaggingRegressor} sortea sus resamples con su propio generador interno,
  controlado por \texttt{random\_state}, que es independiente del que usan las otras
  dos versiones. Sus resamples son \emph{otras} muestras del mismo procedimiento. Por
  lo tanto la comparación correcta es entre distribuciones: los intervalos deben
  solaparse y cubrir $\boldsymbol{\beta}^*$ en proporción similar, pero sus extremos
  difieren en el orden de magnitud de la propia variabilidad Monte Carlo.
\end{description}

\begin{table}[H]
\centering
\caption{Diferencias máximas entre las estimaciones de las distintas versiones.}
\label{tab:equivalencia}
\small
\begin{tabular}{lcl}
\toprule
\textbf{Comparación} & \textbf{Diferencia máxima} & \textbf{Interpretación} \\
\midrule
\texttt{bs\_numpy} vs \texttt{bs\_sklearn}       & $8{,}3 \times 10^{-14}$ & mismas filas $\Rightarrow$ equivalencia numérica \\
\texttt{bs\_numpy} vs \texttt{bs\_numpy (pesos)} & $8{,}5 \times 10^{-14}$ & mismo resample, distinto orden de suma \\
\texttt{bs\_numpy} vs \texttt{bs\_auto}          & $4{,}6 \times 10^{-3}$  & RNG distinto $\Rightarrow$ equivalencia estadística \\
\bottomrule
\end{tabular}
\end{table}

La separación es de once órdenes de magnitud, y confirma nítidamente la distinción
anterior. Las diferencias de $\sim 10^{-14}$ corresponden al error de redondeo de la
aritmética de doble precisión sobre sumas de $100\,000$ términos: son ruido numérico,
no discrepancias algorítmicas. La diferencia de $\sim 10^{-3}$ frente a
\texttt{bs\_auto}, en cambio, es del orden del ancho de los propios intervalos
($\approx 0{,}011$), es decir, exactamente la variabilidad Monte Carlo que cabe esperar
entre dos conjuntos distintos de 48 resamples.

\subsection{Condiciones para que los resultados sean reproducibles}

Para que dos ejecuciones entreguen resultados idénticos deben cumplirse:

\begin{enumerate}
\item \textbf{Semilla fija en la generación de datos.} \texttt{np.random.seed(42)} antes
  de generar $\boldsymbol{\beta}^*$, $X$ y el ruido. Además el \emph{orden} de las
  llamadas al generador debe mantenerse, ya que comparten el mismo estado.

\item \textbf{Semilla fija y determinista en cada tarea.} Cada resample se siembra con
  $42+b$, donde $b$ es el índice del resample, no el del worker. Esta distinción es
  esencial: significa que el resultado \textbf{no depende de $p$}, ni del orden en que
  el planificador asigne las tareas, ni de cuántas le toquen a cada proceso. Si en
  cambio cada worker sembrara desde su propio estado, el resultado cambiaría al variar
  $p$.

\item \textbf{\texttt{random\_state} explícito en \texttt{BaggingRegressor}.} Es el
  único punto donde el sorteo no está bajo nuestro control directo. Sin ese argumento,
  \texttt{scikit-learn} usa el estado global de NumPy en el momento del ajuste y los
  resamples cambian en cada corrida. Esta fue, de hecho, una omisión en nuestra primera
  implementación.

\item \textbf{Mismo orden de reducción en punto flotante.} La suma flotante no es
  asociativa, de modo que $X^\top X$ puede diferir en los últimos bits según cómo BLAS
  particione el trabajo entre threads. Diferencias del orden de $10^{-15}$ son entonces
  esperables entre ejecuciones con distinto número de threads, o entre MKL y OpenBLAS.
  Esto afecta a la reproducibilidad \emph{bit a bit}, no a la validez estadística.
\end{enumerate}

Verificamos empíricamente los puntos 1--3 repitiendo cada versión dos veces dentro de
la misma corrida, a escala completa: tanto \texttt{bs\_numpy} como \texttt{bs\_auto}
arrojaron una diferencia máxima de exactamente $0{,}000 \times 10^{0}$, es decir,
resultados idénticos bit a bit.

% =====================================================================
\section{(d) Creación de procesos por el backend multiprocessing}
\label{sec:d}

\subsection{Qué backend se usa realmente}

\texttt{joblib.Parallel} utiliza por defecto el backend \texttt{loky}, una variante
robustecida de \texttt{multiprocessing} que mantiene un \emph{pool} de procesos
reutilizable entre llamadas sucesivas. A diferencia de \texttt{multiprocessing.Pool},
\texttt{loky} arranca sus workers mediante un mecanismo equivalente a \texttt{spawn}:
lanza un intérprete de Python nuevo que vuelve a importar los módulos necesarios, en
lugar de clonar el proceso padre. Esto lo hace más predecible (evita los problemas
clásicos de \texttt{fork} con threads y con bibliotecas nativas que mantienen estado
interno) a cambio de un costo de arranque mayor.

\subsection{¿Cómo se asigna memoria al crear un proceso?}

Conviene contrastar los dos mecanismos que ofrece Linux, porque el comportamiento de
memoria es muy distinto:

\paragraph{Con \texttt{fork}.} El núcleo no copia la memoria del proceso padre. Duplica
únicamente la tabla de páginas y marca todas las páginas de ambos procesos como
\emph{solo lectura}, compartiendo las páginas físicas. Esto es \textbf{copy-on-write}:
mientras nadie escriba, ambos procesos leen exactamente las mismas páginas físicas y el
costo en memoria es cercano a cero. Cuando alguno intenta escribir, se produce un fallo
de página, el núcleo copia \emph{esa página} (típicamente 4~KiB) y recién ahí se
consume memoria adicional.

Hay una sutileza importante en Python: el intérprete usa conteo de referencias, y el
contador vive en la cabecera de cada objeto. Con solo \emph{leer} un objeto Python se
escribe en su cabecera, lo que dispara la copia de la página que lo contiene. Por eso
el copy-on-write rinde mucho menos de lo esperable para estructuras de datos Python.
Con arreglos NumPy la situación es más favorable: el conteo de referencias toca el
objeto \texttt{ndarray} (unos cientos de bytes), pero \emph{no} el búfer de datos
contiguo, que puede permanecer compartido indefinidamente.

\paragraph{Con \texttt{spawn}/\texttt{loky}.} No hay herencia de memoria en absoluto.
Cada worker es un proceso completamente nuevo, con su propio intérprete, que reimporta
los módulos y recibe los argumentos de la tarea serializados. Todo lo que el worker
necesite debe viajar explícitamente hacia él. El costo de arranque es del orden de
decenas a cientos de milisegundos por proceso, y se paga una sola vez porque
\texttt{loky} reutiliza el pool.

\subsection{¿Qué ocurre con los arreglos \texttt{X} e \texttt{y} al lanzar $p$ procesos?}

Serializar $X$ (241~MB) y enviarlo por un \textit{pipe} a cada uno de los $p$ workers
sería catastrófico: implicaría $p \times 241$~MB de tráfico y de memoria duplicada.
\texttt{joblib} evita esto con un mecanismo específico. Todo arreglo NumPy cuyo tamaño
supere \texttt{max\_nbytes} (por defecto 1~MB) no se serializa: se vuelca una sola vez a
un archivo temporal y se entrega a los workers como un \texttt{numpy.memmap} en modo
solo lectura. Cada worker lo mapea en su espacio de direcciones y el núcleo mantiene una
única copia de las páginas en la caché de páginas, compartida por todos.

En Linux ese archivo se crea en \texttt{/dev/shm} cuando está disponible, es decir, en
un sistema de archivos residente en RAM, de modo que no hay acceso a disco. En nuestra
máquina 2, \texttt{/dev/shm} dispone de 2{,}9~GB.

La consecuencia práctica es que $X$ e $\mathbf{y}$ \textbf{no se replican}: existe una
sola copia física, visible desde los $p$ procesos. Pero esto vale solo para los datos de
entrada. Lo que sí se replica es todo lo que cada tarea \emph{construye}:

\begin{equation}
\underbrace{X}_{\text{241 MB, compartida}}
\quad\longrightarrow\quad
\underbrace{X^{(b)} = X[\texttt{indices}]}_{\text{241 MB, privada por worker}}
\end{equation}

La indexación con un arreglo de índices (\textit{fancy indexing}) siempre materializa un
arreglo nuevo; no existe una vista que represente filas repetidas en orden arbitrario.
De modo que con $p$ procesos el consumo de memoria crece aproximadamente como
$241\,\text{MB} + p \times (241\,\text{MB} + \text{overhead del intérprete})$.

Esta es exactamente la razón por la que la máquina 2, limitada a 6~GB bajo WSL, no puede
ejecutar la versión de \texttt{scikit-learn} con $p = 6$: \texttt{lstsq} necesita, además
de esa copia, espacio de trabajo interno para la SVD, y el sistema operativo termina el
proceso. La variante con pesos de la Sección~\ref{sec:b-iteraciones} ataca precisamente
este punto.

% =====================================================================

% =====================================================================
\section{(e) Oversubscription y threads internos de NumPy}
\label{sec:e}

\subsection{Dos niveles de paralelismo superpuestos}

El esquema tiene paralelismo en dos niveles simultáneos: $p$ procesos lanzados por
\texttt{joblib}, y dentro de cada proceso, los threads que la biblioteca BLAS usa
internamente para las operaciones matriciales. Si ambos niveles se configuran de forma
independiente, el número total de threads activos es $p \times t$, y cuando ese producto
supera el número de cores lógicos se produce \textbf{oversubscription}: más hilos listos
para ejecutar que unidades de ejecución disponibles. El sistema operativo responde
alternando entre ellos, lo que agrega cambios de contexto, invalidaciones de caché y
contención por el ancho de banda de memoria, sin aportar paralelismo real.

\subsection{Threads observados con \texttt{threadpool\_info()}}

En la máquina 1, MKL reporta \textbf{4 threads por defecto} en el proceso padre, sobre 8
cores lógicos. Tomado literalmente, ejecutar $p = 8$ procesos implicaría $8 \times 4 = 32$
threads, una sobresuscripción de $4\times$. Sin embargo, al consultar
\texttt{threadpool\_info()} \emph{desde dentro} de los workers la situación es distinta.

\begin{lstlisting}[language=bash,caption={Salida de \texttt{bs\_numpy.py} con $p=8$ en la máquina 1.},basicstyle=\ttfamily\scriptsize]
--- p=8, t=1 ---
  worker 0: al arrancar [('mkl', 1)] -> dentro de threadpool_limits(1) [('mkl', 1)]
  sin limitar: 8 procesos x 1 threads = 8 sobre 8 cores logicos
--- p=8, t=2 ---
  worker 0: al arrancar [('mkl', 1)] -> dentro de threadpool_limits(2) [('mkl', 2)]
  con t=2: 8 procesos x 2 threads = 16 sobre 8 cores logicos  [OVERSUBSCRIPTION]
--- p=8, t=4 ---
  worker 0: al arrancar [('mkl', 1)] -> dentro de threadpool_limits(4) [('mkl', 4)]
  con t=4: 8 procesos x 4 threads = 32 sobre 8 cores logicos  [OVERSUBSCRIPTION]
\end{lstlisting}

\paragraph{Hallazgo: joblib ya mitiga el problema por defecto.}
Los workers \emph{no} arrancan con los 4 threads del padre, sino con 1. El backend
\texttt{loky} fija automáticamente el número de threads internos a
$\lfloor \text{cores lógicos} / p \rfloor = \lfloor 8/8 \rfloor = 1$, de modo que el
producto $p \times t$ coincide exactamente con el número de cores. \textbf{No hay
oversubscription espontánea.} Ésta aparece únicamente cuando se fuerza $t$ por encima de
ese valor, como en las dos últimas configuraciones.

\subsection{Efecto medido sobre el tiempo}

\begin{table}[H]
\centering
\caption{Tiempo de \texttt{bs\_numpy} con $p = 8$ al forzar distintos valores de $t$.
         Máquina 1, $B = 48$.}
\label{tab:oversub}
\small
\begin{tabular}{crrr}
\toprule
$t$ & \textbf{Threads totales} & \textbf{Tiempo (s)} & \textbf{Penalización} \\
\midrule
1 & $8 \times 1 = 8$  & \textbf{12{,}91} & --- \\
2 & $8 \times 2 = 16$ & 12{,}59 & $-2{,}5\%$ \\
4 & $8 \times 4 = 32$ & 18{,}67 & $+44{,}6\%$ \\
\bottomrule
\end{tabular}
\end{table}

La sobresuscripción moderada ($t = 2$, factor $2\times$) es prácticamente inocua: el
tiempo incluso baja un 2,5\%, dentro del ruido de medición. Esto tiene sentido en un
procesador con \textit{hyperthreading}: los 8 cores lógicos corresponden a 4 físicos, de
modo que parte de los hilos adicionales encuentra unidades de ejecución realmente libres.

La sobresuscripción severa ($t = 4$, factor $4\times$), en cambio, cuesta un
\textbf{44,6\% más de tiempo}. Con 32 hilos compitiendo por 8 cores lógicos, el
planificador alterna constantemente entre ellos; cada cambio de contexto invalida las
cachés privadas y obliga a recargar datos desde memoria, precisamente el recurso que ya
era el cuello de botella.

\paragraph{Cómo se corrige.} Limitando explícitamente los threads internos con
\texttt{threadpool\_limits}, \textbf{dentro de la función que ejecuta cada worker}. Esto
merece énfasis porque nuestra primera implementación lo hacía mal:

\begin{lstlisting}[caption={Patrón incorrecto: el límite no llega a los workers.}]
with threadpool_limits(limits=t, user_api='blas'):   # se aplica al proceso PADRE
    Parallel(n_jobs=p)(delayed(tarea)(...) for b in range(B))
\end{lstlisting}

\texttt{threadpool\_limits} actúa sobre las bibliotecas BLAS \emph{ya cargadas en el
proceso que lo invoca}. Los workers de \texttt{loky} son procesos independientes, que
cargan su propia copia de la biblioteca y no heredan ese límite. El resultado es que el
límite se aplicaba únicamente cuando $p = 1$ (caso en que \texttt{joblib} ejecuta en el
propio proceso padre) y era completamente inefectivo para $p \geq 2$: todas las
mediciones de la grilla $(p,t)$ obtenidas con ese patrón eran inválidas, porque no medían
$t$ en absoluto.

\begin{lstlisting}[caption={Patrón correcto: el límite se aplica en cada worker.}]
def tarea(X, y, semilla, t):
    with threadpool_limits(limits=t, user_api='blas'):   # se aplica al WORKER
        ...
\end{lstlisting}

Una alternativa equivalente es \texttt{Parallel(n\_jobs=p, inner\_max\_num\_threads=t)}.
Optamos por la primera porque permite, además, imprimir \texttt{threadpool\_info()} desde
el interior del worker y así \emph{verificar} que el límite tuvo efecto, en lugar de
suponerlo. Ese es el origen de la salida mostrada más arriba.

\pendiente{\textbf{Faltan las capturas del Monitor del Sistema.} Es la única evidencia
que el script no genera solo, y el enunciado la pide explícitamente. No hay que repetir la
tanda completa: basta correr lo siguiente con el monitor abierto (${\sim}20$~s cada uno) y
capturar la pantalla durante cada ejecución.

\medskip
\texttt{python src/bs\_numpy.py -p 1 -t 1}\\
\texttt{python src/bs\_numpy.py -p 4 -t 1}\\
\texttt{python src/bs\_numpy.py -p 8 -t 1}\\
\texttt{python src/bs\_numpy.py -p 8 -t 4}

\medskip
Guardar la principal como \texttt{resultados/<máquina>/monitor\_sistema.png} y la figura de
abajo la toma automáticamente. Lo que hay que comentar: con $p=1$ se ve un solo core
saturado; con $p=8,\ t=1$ los ocho ocupados de forma pareja; y con $p=8,\ t=4$ los ocho
\emph{igual} de ocupados pero con un $44{,}6\%$ más de tiempo. Ese último contraste es el
punto del ítem: una CPU al 100\,\% no significa que esté rindiendo.}

\begin{figure}[H]
\centering
\figurasiexiste{../resultados/DESKTOP-52S4V2O/monitor_sistema.png}{0.8\textwidth}%
  {Captura del monitor del sistema durante la ejecución.}
\caption{Actividad del sistema durante la ejecución de \texttt{bs\_numpy.py}.}
\label{fig:monitor}
\end{figure}

% =====================================================================
\section{(f) Tiempos de ejecución en función de $p$}
\label{sec:f}

Se ejecutaron las tres versiones con $p \in \{1, 2, \dots, 8\}$, donde $8$ es el número de
cores lógicos de la máquina 1, con $t = 1$ thread interno por proceso.

\begin{table}[H]
\centering
\caption{Tiempos de ejecución $T(p)$ en segundos. Máquina 1, $B = 48$, $t = 1$.
         En negrita el mínimo de cada versión.}
\label{tab:tiempos}
\small
\begin{tabular}{crrr}
\toprule
$p$ & \texttt{bs\_auto} & \texttt{bs\_sklearn} & \texttt{bs\_numpy} \\
\midrule
$T_{\text{serial}}$ & --- & 74{,}68 & 11{,}34 \\
\midrule
1 & \textbf{73{,}05} & 75{,}05 & 15{,}74 \\
2 & 81{,}27 & 64{,}32 & 11{,}48 \\
3 & 90{,}60 & \textbf{61{,}07} & \textbf{9{,}88} \\
4 & 87{,}82 & 61{,}60 & 10{,}58 \\
5 & 73{,}12 & 76{,}31 & 14{,}72 \\
6 & 76{,}19 & 69{,}10 & 12{,}63 \\
7 & 74{,}62 & 72{,}92 & 13{,}06 \\
8 & 86{,}19 & 75{,}77 & 16{,}70 \\
\bottomrule
\end{tabular}
\end{table}

\pendiente{Agregar la tabla equivalente para la máquina 2 y la figura comparativa.}

Tres observaciones:

\begin{enumerate}
\item \texttt{bs\_numpy} es entre $5$ y $7$ veces más rápida que las otras dos para todo
  $p$. La elección del método de ajuste domina por completo sobre la estrategia de
  paralelización.
\item Las tres versiones alcanzan su mínimo muy temprano, en $p = 3$ para
  \texttt{bs\_sklearn} y \texttt{bs\_numpy}, y \emph{empeoran} al seguir agregando
  procesos. En $p = 8$, \texttt{bs\_numpy} tarda más ($16{,}70$~s) que en $p = 1$
  ($15{,}74$~s).
\item La línea base secuencial de \texttt{bs\_numpy} ($11{,}34$~s) es \textbf{más rápida}
  que $T(1)$ con \texttt{joblib} ($15{,}74$~s). No es una contradicción: la corrida
  secuencial deja a MKL usar los 4 threads que trae por defecto, mientras que $T(1)$ se
  midió con $t = 1$. La diferencia cuantifica lo que aporta el paralelismo \emph{interno}
  de BLAS, y es mayor que todo lo que aporta el paralelismo de procesos.
\end{enumerate}

% =====================================================================
\section{(g) Speedup y eficiencia}
\label{sec:g}

\begin{align}
S(p) &= \frac{T(1)}{T(p)}, &
E(p) &= \frac{S(p)}{p}
\end{align}

\subsection{Justificación de la elección de $T(1)$}

Esta elección no es inocua y merece discusión explícita, porque determina por completo la
forma de las curvas.

Al medir $T(1)$ como \texttt{Parallel(n\_jobs=1)}, \texttt{joblib} \emph{no crea ningún
proceso}: ejecuta las tareas secuencialmente en el propio proceso padre. Esa corrida no
paga serialización, ni volcado a \texttt{memmap}, ni arranque de workers. En cambio, ya
con $p = 2$ aparecen de golpe todos esos costos. Es decir, $T(1)$ y $T(p \geq 2)$
\textbf{no miden el mismo experimento}, y usar el primero como referencia mezcla el costo
de paralelizar con el de haber cambiado de régimen de ejecución.

Por eso medimos dos líneas base distintas:

\begin{description}
\item[$T_{\text{serial}}$] --- un ciclo \texttt{for} sin \texttt{joblib}, con BLAS libre
  para usar todos sus threads. Es la implementación secuencial \emph{más rápida
  disponible}, y por lo tanto la referencia honesta: responde a ``¿cuánto gano respecto de
  la mejor forma de resolver esto sin paralelismo de tareas?''.
\item[$T(1)$] --- el backend paralelo con un solo proceso. Aísla el costo de
  \emph{repartir} el trabajo, manteniendo fija la maquinaria de ejecución.
\end{description}

\textbf{Adoptamos $T_{\text{serial}}$} como referencia para $S(p)$ y $E(p)$. Es la elección
más exigente y la más honesta: usar $T(1)$ produciría speedups artificialmente favorables
para \texttt{bs\_numpy}, simplemente porque $T(1)$ está penalizado por $t = 1$. Para
\texttt{bs\_auto} no existe línea base sin \texttt{joblib} --- \texttt{BaggingRegressor}
siempre pasa por él, y \texttt{n\_jobs=1} \emph{es} su modo secuencial ---, de modo que en
esa versión se usa necesariamente $T(1)$, y su columna no es estrictamente comparable con
las otras dos.

\subsection{Resultados}

\begin{table}[H]
\centering
\caption{Speedup $S(p)$ y eficiencia $E(p)$. Máquina 1. Referencia:
         $T_{\text{serial}}$ para \texttt{bs\_sklearn} y \texttt{bs\_numpy}, $T(1)$ para
         \texttt{bs\_auto}.}
\label{tab:speedup}
\small
\begin{tabular}{crrrrrr}
\toprule
& \multicolumn{2}{c}{\texttt{bs\_auto}} & \multicolumn{2}{c}{\texttt{bs\_sklearn}} & \multicolumn{2}{c}{\texttt{bs\_numpy}} \\
$p$ & $S(p)$ & $E(p)$ & $S(p)$ & $E(p)$ & $S(p)$ & $E(p)$ \\
\midrule
1 & 1{,}00 & 1{,}00 & 1{,}00 & 1{,}00 & 0{,}72 & 0{,}72 \\
2 & 0{,}90 & 0{,}45 & 1{,}16 & 0{,}58 & 0{,}99 & 0{,}49 \\
3 & 0{,}81 & 0{,}27 & \textbf{1{,}22} & 0{,}41 & \textbf{1{,}15} & 0{,}38 \\
4 & 0{,}83 & 0{,}21 & 1{,}21 & 0{,}30 & 1{,}07 & 0{,}27 \\
5 & 1{,}00 & 0{,}20 & 0{,}98 & 0{,}20 & 0{,}77 & 0{,}15 \\
6 & 0{,}96 & 0{,}16 & 1{,}08 & 0{,}18 & 0{,}90 & 0{,}15 \\
7 & 0{,}98 & 0{,}14 & 1{,}02 & 0{,}15 & 0{,}87 & 0{,}12 \\
8 & 0{,}85 & 0{,}11 & 0{,}99 & 0{,}12 & 0{,}68 & 0{,}08 \\
\bottomrule
\end{tabular}
\end{table}

\begin{figure}[H]
\centering
\figurasiexiste{../resultados/DESKTOP-52S4V2O/graficos_metricas.png}{\textwidth}%
  {Generar con \texttt{python src/plot\_metrics.py}.}
\caption{Tiempo $T(p)$, speedup $S(p)$ y eficiencia $E(p)$ para las tres versiones,
         comparados con la curva ideal. Máquina 1.}
\label{fig:metricas1}
\end{figure}

\begin{figure}[H]
\centering
\figurasiexiste{../resultados/MAQUINA2/graficos_metricas.png}{\textwidth}%
  {Reemplazar \texttt{MAQUINA2} por la etiqueta real de la segunda máquina.}
\caption{Tiempo $T(p)$, speedup $S(p)$ y eficiencia $E(p)$. Máquina 2.}
\label{fig:metricas2}
\end{figure}

\subsection{Comentario}

El resultado central es contundente y conviene no maquillarlo: \textbf{el speedup máximo
alcanzado es $1{,}22$}, con \texttt{bs\_sklearn} en $p = 3$. \texttt{bs\_numpy} llega a
$1{,}15$, también en $p = 3$. Frente a la curva ideal $S(p) = p$, que en $p = 8$ predeciría
un factor de 8, la ganancia real es esencialmente nula: \textbf{este problema no escala
con el número de procesos en este hardware}.

La eficiencia lo expresa con más crudeza. Cae por debajo de $0{,}5$ ya en $p = 2$ y termina
en $0{,}08$--$0{,}12$ en $p = 8$: con ocho procesos se aprovecha cerca del 10\% de la
capacidad agregada. Dicho de otro modo, siete de cada ocho procesos están efectivamente
esperando.

El comportamiento no-monótono (mínimo en $p = 3$, repunte en $p = 5$, nueva caída en
$p = 8$) se explica por dos efectos superpuestos. Primero, la máquina tiene 4 cores
\emph{físicos} y 8 lógicos: más allá de $p = 4$ los procesos empiezan a compartir unidades
de ejecución vía \textit{hyperthreading}, lo que no duplica el rendimiento en un problema
ya limitado por memoria. Segundo, la granularidad: con $B = 48$ tareas, el reparto es
exacto solo cuando $p$ divide a 48, y para $p = 5$ o $p = 7$ el desbalance impone un techo
adicional.

Vale la pena notar que $S(1) = 0{,}72$ para \texttt{bs\_numpy}: con un solo proceso y
$t = 1$ se es \emph{más lento} que la línea base secuencial, que usa 4 threads de BLAS. Es
la consecuencia directa de haber elegido $T_{\text{serial}}$ como referencia, y deja ver
que en este problema el paralelismo interno de BLAS rinde más que el de procesos.

\pendiente{Repetir el comentario contrastando con la máquina 2, y verificar si el punto
de saturación se mantiene en $p \approx 3$ o se desplaza al disponer de más cores.}

% =====================================================================
\section{(h) Overhead}
\label{sec:h}

El overhead total acumulado sobre los $p$ procesos se define como

\begin{equation}
T_o(p) = p \cdot T(p) - T(1),
\end{equation}

es decir, el trabajo total consumido por los $p$ procesos menos el trabajo útil que
bastaría para resolver el problema secuencialmente. En el caso ideal $T(p) = T(1)/p$ y por
lo tanto $T_o(p) = 0$.

\begin{table}[H]
\centering
\caption{Overhead $T_o(p)$ en segundos-proceso. Máquina 1.}
\label{tab:overhead}
\small
\begin{tabular}{crrr}
\toprule
$p$ & \texttt{bs\_auto} & \texttt{bs\_sklearn} & \texttt{bs\_numpy} \\
\midrule
1 &   0{,}0 &   0{,}4 &   4{,}4 \\
2 &  89{,}5 &  54{,}0 &  11{,}6 \\
3 & 198{,}8 & 108{,}5 &  18{,}3 \\
4 & 278{,}2 & 171{,}7 &  31{,}0 \\
5 & 292{,}5 & 306{,}9 &  62{,}2 \\
6 & 384{,}1 & 339{,}9 &  64{,}5 \\
7 & 449{,}3 & 435{,}8 &  80{,}1 \\
8 & 616{,}5 & 531{,}4 & 122{,}2 \\
\bottomrule
\end{tabular}
\end{table}

\begin{figure}[H]
\centering
\figurasiexiste{../resultados/DESKTOP-52S4V2O/grafico_overhead.png}{0.75\textwidth}%
  {Generar con \texttt{python src/plot\_metrics.py}.}
\caption{Overhead $T_o(p)$ en función del número de procesos. Máquina 1.}
\label{fig:overhead}
\end{figure}

El overhead crece de forma \textbf{aproximadamente lineal en $p$} para las tres versiones,
con pendiente creciente en el tramo final. Para \texttt{bs\_numpy} pasa de $11{,}6$ a
$122{,}2$ segundos-proceso entre $p = 2$ y $p = 8$: un factor de $10{,}5$ mientras $p$ solo
se cuadruplica. Un crecimiento estrictamente lineal indicaría un costo fijo por proceso
(arranque, copia de datos); que crezca más rápido apunta a \emph{contención}, un costo que
cada proceso adicional impone también sobre los ya existentes.

\subsection{Fuentes de overhead en este esquema}

Ordenadas por importancia estimada para este problema:

\begin{enumerate}
\item \textbf{Contención por ancho de banda de memoria.} Es, con diferencia, la fuente
  dominante, y la evidencia la respalda por tres vías independientes: el speedup satura en
  $1{,}22$; el overhead crece superlinealmente; y la variante con pesos
  (Sección~\ref{sec:b-iteraciones}) no mejoró pese a optimizar el patrón de acceso. Cada
  tarea recorre $X$ (241~MB) para construir $X^{(b)}$ y luego lo recorre otra vez para
  calcular $X^{(b)\top}X^{(b)}$. Los cores comparten un único bus hacia la RAM: agregar
  procesos agrega unidades de cómputo, pero \emph{no} agrega ancho de banda.

\item \textbf{Materialización de $X^{(b)}$.} Cada worker asigna y llena 241~MB por
  resample. Con $B = 48$ resamples eso son $\sim 11{,}6$~GB de tráfico de memoria por
  corrida, únicamente para construir las muestras.

\item \textbf{Presión sobre la caché compartida.} Con $p$ procesos coexisten $p$ copias
  privadas de 241~MB, que se expulsan mutuamente de la caché L3 compartida. Esto degrada
  el rendimiento de todos los procesos a la vez, y explica por qué el costo es
  superlineal.

\item \textbf{Desbalance de carga (granularidad).} Con $B = 48$ tareas el reparto es
  perfecto solo cuando $p$ divide a 48. Para $p = 5$ algunos workers reciben 10 tareas y
  otros 9, y el total lo fija el más cargado: la eficiencia máxima alcanzable cae a
  $48/(10 \cdot 5) = 96\%$. Para $p = 7$ el techo es $48/(7 \cdot 7) = 98\%$. Este efecto
  contribuye a la irregularidad de las curvas, aunque es menor frente a la contención.

\item \textbf{Creación de procesos y transferencia de datos.} Arrancar los workers de
  \texttt{loky} (intérpretes nuevos que reimportan NumPy y scikit-learn) y volcar $X$ a
  \texttt{memmap} tiene un costo fijo, del orden de cientos de milisegundos. Es
  \emph{constante} respecto de $B$, y por eso pesa mucho más en \texttt{bs\_numpy}
  ($\sim 10$~s totales) que en las versiones de scikit-learn ($\sim 70$~s).

\item \textbf{Oversubscription de threads.} Cuantificada en la Sección~\ref{sec:e}: un
  $44{,}6\%$ de penalización con $t = 4$. Está controlada en estas mediciones ($t = 1$),
  pero sería la fuente dominante si no se gestionara.

\item \textbf{Hyperthreading.} Más allá de $p = 4$ (cores físicos) los procesos comparten
  unidades de ejecución. En un problema limitado por memoria, el segundo hilo de cada core
  aporta poco y compite por la misma caché.
\end{enumerate}

% =====================================================================
\section{(i) Grilla de procesos y threads}
\label{sec:i}

Usando \texttt{bs\_numpy}, la versión más eficiente del ítem (b), se exploraron
combinaciones $(p, t)$ sujetas a $p \cdot t \leq p_{\max}$, de modo de no exceder nunca el
número de cores lógicos. Las combinaciones que violan esa restricción se excluyeron.

\begin{table}[H]
\centering
\caption{Tiempo $T(p,t)$ en segundos. Máquina 1, $p_{\max} = 8$, $B = 48$.
         Las celdas marcadas con --- violan $p \cdot t \leq 8$.}
\label{tab:grid}
\small
\begin{tabular}{crrrr}
\toprule
& \multicolumn{4}{c}{\textbf{Procesos $p$}} \\
\textbf{Threads $t$} & $1$ & $2$ & $4$ & $8$ \\
\midrule
1 & 17{,}31 & 13{,}50 & \textbf{11{,}10} & 12{,}70 \\
2 & 13{,}68 & 11{,}14 & 11{,}75 & --- \\
4 & 11{,}49 & 11{,}41 & --- & --- \\
8 & 11{,}70 & --- & --- & --- \\
\bottomrule
\end{tabular}
\end{table}

\begin{figure}[H]
\centering
\figurasiexiste{../resultados/DESKTOP-52S4V2O/grid_heatmap.png}{0.75\textwidth}%
  {Generar con \texttt{python src/run\_grid\_experiment.py}.}
\caption{Tiempo de ejecución $T(p,t)$ para cada combinación válida. Máquina 1.}
\label{fig:grid1}
\end{figure}

\begin{figure}[H]
\centering
\figurasiexiste{../resultados/MAQUINA2/grid_heatmap.png}{0.75\textwidth}%
  {Reemplazar \texttt{MAQUINA2} por la etiqueta real de la segunda máquina.}
\caption{Tiempo de ejecución $T(p,t)$. Máquina 2.}
\label{fig:grid2}
\end{figure}

\subsection{Conclusión}

La mejor combinación es $\boldsymbol{(p, t) = (4, 1)}$, con $11{,}10$~s. Pero el hallazgo
más interesante es \textbf{lo plana que resulta la superficie}: todas las configuraciones
que saturan los 8 cores lógicos --- $(4,1)$, $(2,2)$, $(1,4)$, $(1,8)$, $(2,4)$ ---
quedan entre $11{,}10$ y $11{,}75$~s, un rango de apenas $6\%$.

Esto refuerza la interpretación de la Sección~\ref{sec:h}: si el recurso escaso fuera la
capacidad de cómputo, la forma de repartirla entre procesos y threads importaría. Como el
recurso escaso es el \emph{ancho de banda de memoria}, y ambos niveles de paralelismo
compiten por el mismo bus, da casi igual cómo se distribuya el trabajo mientras se ocupen
todos los cores.

Las configuraciones claramente malas son las que \emph{no} saturan la máquina: $(1,1)$ con
$17{,}31$~s es la peor de todas, usando un solo core. Y $(8,1)$ con $12{,}70$~s es un
$14\%$ peor que $(4,1)$, consistente con que los 8 cores lógicos son solo 4 físicos.

\paragraph{Compromiso entre los dos niveles.} El paralelismo de tareas (aumentar $p$) no
requiere comunicación entre resamples, pero multiplica el uso de memoria: cada proceso
mantiene su propia copia de $X^{(b)}$. El paralelismo de threads (aumentar $t$) comparte
los datos dentro del proceso, pero escala mal en operaciones limitadas por ancho de banda.
En esta máquina el óptimo cae en un punto intermedio con predominio de procesos, aunque la
diferencia con las alternativas es marginal.

\pendiente{Comparar con la grilla de la máquina 2 y verificar si el óptimo se mantiene en
$t$ bajo o se desplaza al haber más cores disponibles.}

% =====================================================================
\section{(j) Comparación entre los dos computadores}
\label{sec:j}

\pendiente{Sección pendiente de la tanda de la máquina 2. La columna de la máquina 1 ya
está completa; falta la segunda y el comentario.

\medskip
\textbf{Estructura sugerida.}

\emph{1. Tabla comparativa:} $T_{\text{serial}}$, $T(p_{\max})$, mejor tiempo y a qué $p$,
speedup máximo, eficiencia en $p_{\max}$, y mejor combinación $(p,t)$.

\emph{2. Diferencias esperables y su justificación:}
\begin{itemize}
  \item \emph{Número de cores.} No debería traducirse en un speedup proporcional,
    precisamente porque el problema está limitado por ancho de banda: más cores no
    implican más ancho de banda hacia la RAM. Verificar si el punto de saturación
    ($p \approx 3$ en la máquina 1) se desplaza o se mantiene.
  \item \emph{Relación cores físicos / lógicos.} En la máquina 1 el quiebre en $p = 4$
    coincide con el número de cores físicos. Comprobar si el mismo patrón aparece.
  \item \emph{Biblioteca BLAS y su default de threads.} Comparar los tiempos con $p=1$,
    que es donde el efecto de BLAS se aísla mejor.
  \item \emph{Jerarquía de memoria.} Tamaño de L3 y número de canales de memoria explican
    buena parte del punto de saturación.
  \item \emph{Memoria disponible.} Con poca RAM, las versiones de scikit-learn no llegan a
    $p$ alto: \texttt{lstsq} usa $\sim 1$~GB por proceso.
\end{itemize}

\emph{3. Anécdota de portabilidad.} \texttt{run\_grid\_experiment.py} contenía una
f-string con una barra invertida, admitida en Python 3.12+ (PEP 701) pero que \emph{no
compila} en 3.10. El código corría en una máquina y fallaba al importarse en la otra. Las
diferencias entre entornos no se limitan al rendimiento.}

% =====================================================================
\section{Conclusiones}
\label{sec:conclusiones}

\begin{enumerate}
\item \textbf{La optimización algorítmica superó ampliamente a la paralelización.}
  Reemplazar la descomposición SVD de \texttt{lstsq} por la resolución directa de la
  ecuación normal produjo una mejora de $\sim 6\times$ ($75$~s a $11$~s). Todo el
  paralelismo de procesos, en cambio, aportó a lo sumo un factor $1{,}22$. Antes de
  paralelizar conviene revisar el algoritmo.

\item \textbf{El problema está limitado por ancho de banda de memoria, no por cómputo.}
  Tres evidencias independientes convergen: el speedup satura en $p \approx 3$ con un
  factor de apenas $1{,}22$; el overhead crece superlinealmente en $p$; y la superficie
  $T(p,t)$ es esencialmente plana ($6\%$ de rango) entre todas las configuraciones que
  ocupan los 8 cores. Con $k = 301$, la razón entre operaciones aritméticas y bytes
  movidos es demasiado baja para que agregar procesos ayude.

\item \textbf{No toda micro-optimización paga.} La variante que reemplaza el
  \textit{fancy indexing} por pesos multinomiales, matemáticamente equivalente y con mejor
  patrón de acceso, resultó un $24\%$ \emph{más lenta}. Lo que domina es el volumen de
  datos movido, no el orden en que se accede.

\item \textbf{Las tres implementaciones son correctas y equivalentes.}
  \texttt{bs\_numpy} y \texttt{bs\_sklearn} coinciden hasta $8 \times 10^{-14}$ --- ruido
  de redondeo --- porque comparten semilla y remuestrean las mismas filas.
  \texttt{bs\_auto} difiere en $4{,}6 \times 10^{-3}$, del orden del propio ancho de los
  intervalos, porque usa su generador interno: su equivalencia es estadística, no
  numérica. Las cuatro cubren $\boldsymbol{\beta}^*$ en torno al $91\%$ de los
  coeficientes, consistente con el $95\%$ nominal dado que $B = 48$ deja muy pocas
  observaciones en cada cola.

\item \textbf{Medir correctamente es parte del problema.} Dos decisiones metodológicas
  cambiaron las conclusiones por completo: aplicar \texttt{threadpool\_limits} dentro del
  worker y no en el proceso padre (sin lo cual la grilla del ítem (i) no medía nada), y
  elegir una línea base secuencial honesta en lugar de $T(1)$ con \texttt{joblib} (que
  infla los speedups). Ambas fueron errores de nuestra primera implementación, detectados
  al verificar que las mediciones dijeran lo que creíamos.
\end{enumerate}

\pendiente{Revisar y ampliar con las conclusiones que surjan de la comparación entre
máquinas (ítem j).}

% =====================================================================
\section*{Declaración de uso de inteligencia artificial}
\addcontentsline{toc}{section}{Declaración de uso de inteligencia artificial}

\pendiente{Revisar y ajustar esta declaración para que refleje con exactitud el uso que el
grupo le dio a herramientas de IA. El texto siguiente es una base.}

En la realización de esta tarea se utilizó el asistente de programación \textbf{Claude
Code} (Anthropic), con los siguientes propósitos:

\begin{itemize}
\item \textbf{Revisión de código.} Identificación de errores en la implementación
  original, en particular: la aplicación incorrecta de \texttt{threadpool\_limits} en el
  proceso padre en lugar de en los workers (Sección~\ref{sec:e}); la ausencia de
  \texttt{random\_state} en \texttt{BaggingRegressor}, que impedía la reproducibilidad
  (Sección~\ref{sec:c}); el incumplimiento de la restricción $p \cdot t \leq p_{\max}$ en
  la grilla del ítem (i); y la ausencia de una línea base secuencial para el cálculo del
  speedup (Sección~\ref{sec:g}).
\item \textbf{Refactorización.} Parametrización de los scripts, separación de resultados
  por máquina y automatización de la tanda de experimentos.
\item \textbf{Redacción.} Estructuración y redacción de este informe a partir de los
  resultados experimentales obtenidos por el grupo.
\end{itemize}

El diseño experimental, la ejecución de las mediciones en ambos computadores y la
interpretación de los resultados fueron realizados por los integrantes del grupo. Todo el
código fue revisado y verificado por el grupo antes de su uso.

\end{document}
